\chapter{Example: photon-assisted tunnelling into a driven junction}\label{ch:ex-pat}

A weak tunnelling contact, such as a scanning-tunnelling tip or a normal lead, adds or
removes single electrons and thereby measures transitions between the even and odd subgap
states. Under a microwave gate drive each tunnelling line acquires sidebands at multiples of
the drive frequency: photon-assisted tunnelling. It was described by Tien and Gordon for
superconducting tunnel junctions~\cite{tien1963multiphoton}, observed in quantum
dots~\cite{kouwenhoven1994observation,oosterkamp1998microwave}, and recently used to probe
Yu--Shiba--Rusinov states with a scanning tunnelling
microscope~\cite{kot2020microwaveassisted,peters2020resonant,acerogonzalez2020photonassisted}.

\section{Tien--Gordon and its generalization}

For the decoupled dot, \cref{ch:ex-dot}, the sideband weights are Bessel functions,
$J_q^2(\Delta n\,A/\Omega)$, with $\Delta n$ the charge difference between the final and the
initial state. This result is exact, and the tests of \pkg{} check that the code reproduces it
to $10^{-10}$. In the junction the subgap states have non-integer charges, $n_S=\q{pat.n_S}$
and $n_D=\q{pat.n_D}$ for the bath used here, so the simple Tien--Gordon formula with
$\Delta n=\q{pat.dn}$ provides a natural first approximation. It neglects two effects. The transition energy
$E_D(\eps)-E_S(\eps)$ depends nonlinearly on the detuning, and so does the tunnelling matrix
element $\matel D{d^\dagger_\uparrow}S$. For a slow drive both are captured by an adiabatic
generalization, in which the sideband amplitudes are the Fourier coefficients of
\begin{equation}\label{eq:adiabatic-tg}
  F(t)=\mathcal M\bigl(\eps(t)\bigr)\exp\Bigl[\ii\int_0^t\bigl(\Delta E(\eps(t'))-\overline{\Delta E}\bigr)\dd t'\Bigr],
\end{equation}
with $\mathcal M(\eps)=\matel{D(\eps)}{d^\dagger_\uparrow}{S(\eps)}$ and $\Delta E=E_D-E_S$ from static
calculations along the drive.

\section{Results}

In \cref{fig:pat} the gate of a junction with $\q{pat.pairs}$ bath level pairs per lead is
driven at $\Omega=\q{pat.omega}$, far below the excitation gaps within each parity sector
($\q{pat.gap_singlets}$ for the singlets, $\q{pat.gap_doublet}$ above the doublet), and the
addition spectrum of $d^\dagger_\uparrow$ is evaluated from the dressed singlet. Panel (a)
shows the sidebands spreading out as the amplitude grows. Panel (b) compares the weights of
the lines into the dressed doublet with the two formulas. The Floquet results follow the
adiabatic generalization \eqref{eq:adiabatic-tg} closely, while the simple Tien--Gordon expression reproduces
the overall envelope with symmetric weights for $q=+1$ and $q=-1$. The adiabatic expression captures
the asymmetry through the detuning-dependent energy and matrix element. For $A=\q{pat.range_amplitude}$, $\Delta E$ varies from $\q{pat.gap_min}$ to
$\q{pat.gap_max}$ and $\abs{\mathcal M}$ from $\q{pat.element_min}$ to $\q{pat.element_max}$ over the
drive period. Evaluating
\eqref{eq:adiabatic-tg} with a constant $\mathcal M$, or with a linearized $\Delta E$, shows that the two
dependences produce asymmetries of opposite sign; the observed asymmetry is their net
effect. The Floquet calculation retains nonadiabatic dynamics and extends the
treatment to frequencies $\Omega$ approaching the internal gaps, beyond the regime of \eqref{eq:adiabatic-tg}.

\begin{figure}[htb]
\centering
\manualfigure{pat}
\caption{Photon-assisted tunnelling: gate drive at $\Omega=\q{pat.omega}$ on a junction with
$\q{pat.pairs}$ bath level pairs per lead ($U=\q{pat.u}$, $\Gamma=\q{pat.gamma}$,
$\eps=\q{pat.epsilon}$, $\phi=\q{pat.phi}$). (a) Addition spectrum of $d^\dagger_\uparrow$
from the dressed singlet for several amplitudes, offset vertically (Lorentzian half width
$\q{pat.width}$).
(b) Weights of the sidebands into the dressed doublet, normalized to the static
strength. Markers: Floquet; solid ($q\ge0$) and dashed ($q<0$) lines: adiabatic generalized
Tien--Gordon \eqref{eq:adiabatic-tg}; dotted lines: plain $J_q^2(\Delta n\,A/\Omega)$.}
\label{fig:pat}
\end{figure}

\begin{caution}
Microwaves applied to a tunnelling experiment also modulate the probe contact itself, which
produces replicas of every line independently of any dressing of the junction. Separating
the two requires the drive to act on the junction alone, or a model of the probe
circuit~\cite{haxell2023microwaveinduced}.
\end{caution}

\tutorial{09_photon_assisted_tunnelling}
